\documentclass[../../main.tex]{subfiles}
\begin{document}

\begin{center}
  {\fontsize{48pt}{58pt}\selectfont \textbf{2011}\normalsize\textbf{年度}}
\end{center}

\begin{table}[htb]
  \centering
  \caption{出題テーマと難易度}
  \begin{tabular}{|c|p{9cm}|c|}
    \hline
    \textbf{問題番号} & \textbf{テーマ} & \textbf{難易度} \\
    \hline
    第1問 & カードの小さい番号の一致確率・定積分 & \\
    \hline
    第2問 & 行列で表す一次変換と面積 & \\
    \hline
    第3問 & 直線と絶対値曲線で囲む面積 & \\
    \hline
    第4問 & 積の不等式と数学的帰納法 & \\
    \hline
    第5問 & 球面と平面の交円上での積の範囲 & \\
    \hline
    第6問 & 四面体の外接球の存在 & \\
    \hline
  \end{tabular}
\end{table}

\begin{table}[htb]
  \centering
  \caption{解答}
  \renewcommand{\arraystretch}{1.4}
  \begin{tabular}{|c|c|p{9cm}|}
    \hline
    \textbf{問題番号} & \textbf{小問} & \textbf{答え} \\
    \hline
    第1問 & (1) & $\dfrac{17}{108}$ \\
    \cline{2-3}
    & (2) & $\dfrac{1}{6}+\dfrac{\sqrt{2}}{4}\left(\dfrac{\pi}{4}+\dfrac{1}{3}\right)$ \\
    \hline
    第2問 & -- &  \\
    \hline
    第3問 & -- & $\dfrac{208}{27}$ \\
    \hline
    第4問 & -- & 指定された積の不等式が成立（証明問題） \\
    \hline
    第5問 & -- & $\dfrac{50}{27}\le xyz\le2$ \\
    \hline
    第6問 & -- & 4頂点を通る球面が存在（証明問題） \\
    \hline
  \end{tabular}
\end{table}

\paragraph{試験形態}
6問．試験時間は要確認．

\end{document}
